\documentclass[11pt]{article}
\usepackage{fontspec}
\setmainfont{Times New Roman}
\setsansfont{Arial}
\setmonofont{Consolas}
\usepackage{amsmath,amssymb,mathtools}
\usepackage{amsthm}
\usepackage{geometry}
\usepackage{microtype}
\geometry{a4paper,margin=2.05cm}
\setlength{\parindent}{1.1em}
\setlength{\parskip}{0pt}
\newcommand{\foreign}[1]{#1}
\theoremstyle{plain}
\newtheorem{satz}{Теорема}
\newtheorem{zusatz}{Додаток}
\emergencystretch=220em
\allowdisplaybreaks
\sloppy
\hbadness=10000
\vbadness=10000
\hfuzz=5pt
\vfuzz=12pt
\raggedbottom
\begin{document}

\newpage
\part*{Частина III\hspace{1em}Стаття Капферера --- Нетер}
\addcontentsline{toc}{part}{Частина III: Стаття Капферера --- Нетер}

\begin{center}
\Large\textbf{Необхідні й достатні умови кратності}\\[0.3em]
\Large\textbf{для Нетерової основної теореми}\\[0.5em]
\large\textbf{алгебраїчних функцій}\\[1em]
\normalsize\textit{Г. Капферер}\\[0.5em]
\small\textit{(з додатком, спільно з Е. Нетер)}\\[1em]
\normalsize \foreign{Math.\ Ann.}\ \textbf{97} (1927), S.~559--567
\end{center}

\vspace{1em}

Нижче подано посилення Нетерової основної теореми у формі
необхідних і достатніх умов кратності\textsuperscript{1)}. Комплекс
фактів Нетерової основної теореми я, спираючись на
ідеально-теоретичне розуміння цієї теореми\textsuperscript{2)}, але
не користуючись тут і далі ідеально-теоретичними поняттями та
теоремами, зводжу до таких двох окремих тверджень\textsuperscript{3)}:

\smallskip
\textbf{Теорема Ia.} Для кожного спільного нуля $x = a$, $y = b$
двох взаємно простих поліномів $\varphi$, $\psi$ від $x$, $y$ з
довільними комплексними коефіцієнтами існує найменше ціле
раціональне додатне число $\varrho$ таке, що модуль
$(\varphi, \psi, \mathfrak{p}^{\varrho})$ означає точно ту саму
сукупність поліномів, що й модуль
$(\varphi, \psi, \mathfrak{p}^{\sigma})$, де
$\mathfrak{p} = (x-a, y-b)$, а $\sigma$ означає довільне ціле
раціональне число, більше за $\varrho$\textsuperscript{4)}.

\smallskip
\textbf{Теорема Ib.} Якщо цю поняттєво однозначно визначену
сукупність поліномів позначити через $\mathfrak{q}$, відповідно
через $\mathfrak{q}_{\nu}$, коли йдеться про $\nu$-й із $s$ різних
нулів, тобто про $x = a_{\nu}$, $y = b_{\nu}$, то має місце таке:

Необхідно й достатньо для того, щоб заданий поліном $K$ від $x$, $y$
мав властивість подільності $K \equiv 0(\varphi, \psi)$, щоб
одночасно виконувалися такі $s$ конгруенцій
\[
K \equiv 0(\mathfrak{q}_{1}), \quad K \equiv 0(\mathfrak{q}_{2}), \quad\ldots,\quad
K \equiv 0(\mathfrak{q}_{s}).
\]

Історично, порівняймо вступ до основоположної праці Макса Нетер
про цей предмет\textsuperscript{5)}, основна теорема постала зі
спроб розв'язати таку фундаментальну задачу: для кожного заданого
полінома $K$ від $x$, $y$ треба вміти встановити, чи він ділиться
на $(\varphi, \psi)$, тобто також чи існують два інші поліноми
$A$ і $B$ від $x$, $y$ такі, що рівність
\[
K = A\varphi + B\psi
\]
стає тотожністю відносно $x$, $y$.

Цю задачу теорема Нетер не розв'язує; радше її треба розуміти лише
як, хоч і фундаментальне, зведення\textsuperscript{6)} поставленої
задачі до простішої, а саме до питання про необхідні й достатні
умови існування окремих конгруенцій $K \equiv 0(\mathfrak{q})$.
Коли ж має місце $K \equiv 0(\mathfrak{q})$? Лише після відповіді
на це останнє питання розв'язано і початкову постановку задачі
М. Нетер. Саме це є метою цієї статті, і ця мета справді
досягається: вказуються скінченно багато умов кратності, які
практично також неважко перевіряти і які є необхідними й
достатніми для $K \equiv 0(\mathfrak{q})$. Одна достатня умова для
$K \equiv 0(\mathfrak{q})$, яка однак не є водночас необхідною
умовою, відома і часто застосовується та цитується в літературі,
наприклад у такій формі:

\smallskip
\textit{Для існування конгруенції $K \equiv 0(\varphi, \psi)$
достатньо, щоб кожна точка перетину $\varphi = 0$, $\psi = 0$
входила до полінома $K$ з ``досить'' високою кратністю.}

\smallskip
Ця теорема міститься в теоремі I як спеціальний випадок; адже
$\varrho$-кратна точка в $K$ рівнозначна $K \equiv 0(\mathfrak{p}^{\varrho})$;
тоді тим більше $K \equiv 0(\varphi, \psi, \mathfrak{p}^{\varrho})$,
тобто також $K \equiv 0(\mathfrak{q})$. Але обернене твердження не
має місця.

\vspace{0.5em}
\noindent\rule{\textwidth}{0.4pt}
\vspace{0.3em}

{\footnotesize
\textsuperscript{1)} Докладний виклад я відсилаю до моєї праці з тією самою
назвою у спеціальному випуску звітів Гайдельберзької академії 1927 року
``Beiträge zur Algebra'', ініційованому А.\ Леві, Фрайбург-ім-Брайсгау, який
містить праці різних авторів.

\textsuperscript{2)} Пор., наприклад, B.\ L.\ van der Waerden, ``Zur Nullstellentheorie der Polynomideale''.
\foreign{Math.\ Annalen} 96 (1926), S.~203.

\textsuperscript{3)} У \S~1 цитованої праці я довів обидва твердження окремо й
безпосередньо, тобто без ідеально-теоретичних понять.

\textsuperscript{4)} Літери $\mathfrak{p}$ і $\mathfrak{q}$ вибрано для того, щоб нагадати про
змістову відповідність позначуваного поняттям простого ідеалу
$\mathfrak{p}$ і первинного ідеалу $\mathfrak{q}$, усталеним у новішій літературі.

\textsuperscript{5)} Max Noether, ``Über einen Satz aus der Theorie der algebraischen Funktionen'',
\foreign{Math.\ Annalen} 6 (1872).

\textsuperscript{6)} Згадане зведення приблизно порівнянне з тим, яке кожен знає з
елементарної теорії чисел: властивість подільності натурального числа
$A$ на таке саме $B$ еквівалентна властивості подільності $A$ на
$p_{1}^{e_{1}}$, на $p_{2}^{e_{2}}$, \ldots, на $p_{s}^{e_{s}}$, де
$p_{1}, p_{2}, \ldots, p_{s}$ означають різні прості дільники
$B = p_{1}^{e_{1}}p_{2}^{e_{2}} \cdots p_{s}^{e_{s}}$. Саме поняття
``первинний ідеал'', у тексті позначене через $\mathfrak{q}$, треба
розглядати як аналог степеня простого числа. [Пор. E.\ Noether,
``Idealtheorie in Ringbereichen'', \foreign{Math.\ Annalen} 83 (1921), вступ.]
}

\newpage
\section*{Головна теорема}

Оголошений розв'язок задачі можна коротко охарактеризувати так:
кожному окремому $\mathfrak{q}$ можна поставити у відповідність
систему з $t$ додатних цілих раціональних чисел
$(g_{1}), (g_{2}), \ldots, (g_{t})$, а модулю $(K, \mathfrak{q})$ таку
саму кількість цілих раціональних додатних або від'ємних чисел
$(K_{1}), (K_{2}), \ldots, (K_{t})$, так що має місце така
\textbf{головна теорема}:

\begin{satz}[Головна теорема]\label{satz:kap-II}
Необхідно й достатньо для $K \equiv 0(\mathfrak{q})$, щоб одночасно
виконувалися $t$ умов
\[
(K_{1}) \geq (g_{1}), \quad (K_{2}) \geq (g_{2}), \quad\ldots,\quad (K_{t}) \geq (g_{t}).
\]
При цьому
\[
(g_{1}) + (g_{2}) + \cdots + (g_{t}) = \mu,
\]
де $\mu$ означає відповідну до $\mathfrak{q}$ кратність перетину,
визначену результантом.
\end{satz}

У доведенні, без обмеження загальності, вважатимемо $x = 0$,
$y = 0$ нулем пари $(\varphi, \psi)$, що відповідає
$\mathfrak{q}$. Далі взагалі, якщо $A(x,y)$ є будь-яким
поліномом\textsuperscript{7)} від $x$, $y$, символ $(A)$ означатиме
число, яке показує, скільки разів $y$ входить як множник в $A(0,y)$;
якщо ж $A(0,y)$ тотожно зникає, то $(A)$ покладається рівним
$\infty$. Числа $(g_{1}), (g_{2}), \ldots, (g_{t})$, поставлені у
відповідність $\mathfrak{q}$, визначаються такою системою рівнянь,
яка становить методичну основу процедури:
\begin{equation}\label{eq:kap-1}
\begin{aligned}
f_{1} \cdot g_{1}(0,y):y^{(g_{1})} - g_{1} \cdot f_{1}(0,y):y^{(g_{1})} &= x \cdot h_{1},\\
f_{2} \cdot g_{2}(0,y):y^{(g_{2})} - g_{2} \cdot f_{2}(0,y):y^{(g_{2})} &= x \cdot h_{2},\\
&\vdots \qquad\qquad\vdots\qquad\qquad\vdots\\
f_{r} \cdot g_{r}(0,y):y^{(g_{r})} - g_{r} \cdot f_{r}(0,y):y^{(g_{r})} &= x \cdot h_{r}.
\end{aligned}
\end{equation}
Вона будується лише з заданої пари поліномів $\varphi$, $\psi$:
поліноми $f_{1}, g_{1}$ визначаються як тотожні з $\varphi$, $\psi$,
можливо після перестановки; порядок має бути таким, щоб
$(f_{1}) \geq (g_{1})$. За цієї фіксації, і лише тоді, $h_{1}$ буде
ціло-раціональним.

Поліноми $f_{2}, g_{2}$ визначаються як тотожні з $h_{1}, g_{1}$ за
додаткової умови $(f_{2}) \geq (g_{2})$. Загалом пара поліномів
$f_{r}, g_{r}$ визначається як тотожна з $h_{r-1}, g_{r-1}$ за
додаткової умови $(f_{r}) \geq (g_{r})$.

Систему рівнянь \eqref{eq:kap-1} справді можна продовжувати без кінця
у пояснений спосіб; але головна теорема вимагає побудувати лише
перші $t$ рівнянь системи, де $t$ задається таким фактом:

\smallskip
\textbf{Допоміжна теорема I.} Існує натуральне число $t$ [і саме
$t \leq \mu$, де під $\mu$ розуміється кратність перетину, з якою
точка $x = 0$, $y = 0$ входить до $\varphi = 0$, $\psi = 0$] таке, що
перші $t$ чисел, визначених через \eqref{eq:kap-1}, а саме
$(g_{1}), (g_{2}), \ldots, (g_{t})$, усі більші за нуль, тоді як
водночас $(g_{t+1})$ і всі наступні дорівнюють нулю.

Справді, з урахуванням теореми про добуток результантів кратність
перетину в точці $x = 0$, $y = 0$ для $f_{r} = 0$, $g_{r} = 0$
дорівнює
\[
\mu - (g_{1}) - (g_{2}) - \cdots - (g_{r}),
\]
звідки випливає обрив після щонайбільше $\mu$ кроків. З цього
одночасно випливає факт
\begin{equation}\label{eq:kap-2}
\mu = (g_{1}) + (g_{2}) + \cdots + (g_{t}),
\end{equation}
формула, примітна тим, що вона загалом дає раціональний і практично
легко здійсненний метод визначення кратності перетину, щойно відомий
відповідний нуль.

Коли числа $(g_{i})$ уже відомі, числа $(K_{i})$ визначаються такою
системою рівнянь, відповідною до \eqref{eq:kap-1}:
\begin{equation}\label{eq:kap-3}
\begin{aligned}
K_{1} \cdot g_{1}(0,y):y^{(g_{1})} - g_{1} \cdot K_{1}(0,y):y^{(g_{1})} &= x \cdot K_{2},\\
K_{2} \cdot g_{2}(0,y):y^{(g_{2})} - g_{2} \cdot K_{2}(0,y):y^{(g_{2})} &= x \cdot K_{3},\\
&\vdots \qquad\qquad\vdots\qquad\qquad\vdots\\
K_{r} \cdot g_{r}(0,y):y^{(g_{r})} - g_{r} \cdot K_{r}(0,y):y^{(g_{r})} &= x \cdot K_{r+1}.
\end{aligned}
\end{equation}
$K_{1}$ визначається як тотожний з $K$. Очевидно, $K_{2}$ буде
ціло-раціональним тоді й лише тоді, коли $(K_{1}) \geq (g_{1})$.
Якщо останнє виконано, далі випливає: $K_{3}$ буде
ціло-раціональним тоді й лише тоді, коли $(K_{2}) \geq (g_{2})$.
Загалом, якщо $K_{r}$ уже ціло-раціональний, $K_{r+1}$ буде
ціло-раціональним тоді й лише тоді, коли $(K_{r}) \geq (g_{r})$.

Доведення головної теореми (теореми~\ref{satz:kap-II}) тепер
спирається на ще одну

\smallskip
\textbf{Допоміжну теорему II.} Якщо покласти
$(f_{i}, g_{i}, \mathfrak{p}^{e})$ рівним $\mathfrak{q}^{(i)}$, то для
кожного $i$ з ряду цілих чисел $1, 2, \ldots$ \textit{in inf.} має
місце твердження:

Для конгруенції
\[
K_{i} \equiv 0(\mathfrak{q}^{(i)})
\]
необхідним і достатнім є одночасне виконання двох умов:
\[
(K_{i}) \geq (g_{i}) \qquad \text{і} \qquad
K_{i+1} \equiv 0(\mathfrak{q}^{(i+1)}).
\]

Доведення\textsuperscript{8)} допоміжної теореми дістається
комбінацією систем рівнянь \eqref{eq:kap-1} і \eqref{eq:kap-3}. Якщо
врахувати, що $g_{t+1}$ за допоміжною теоремою I вже не зникає в
нулі, тобто що модуль $\mathfrak{q}^{(t+1)}$ складається з усіх
поліномів (є одиничним ідеалом), і отже
$K_{t+1} \equiv 0(\mathfrak{q}^{(t+1)})$ стає тривіальним, то
$t$-кратне застосування допоміжної теореми II дає доведення
головної теореми.

\smallskip
Без доведень, вони містяться в моїй цитованій вище праці, додам ще
два додатки до головної теореми:

\begin{zusatz}[I]
Число $t$ головної теореми тотожне з найменшим\textsuperscript{9)}
додатним числом $\nu$ такого типу, що $x^{\nu} \equiv 0(\mathfrak{q})$.
\end{zusatz}

\begin{zusatz}[II]
Системи рівнянь \eqref{eq:kap-1} і \eqref{eq:kap-3}, а також ті, що
виникають із них перестановкою $x$ і $y$, за придатної комбінації
дають раціональний метод точного числового визначення ``належного''
до $\mathfrak{q}$ показника $\varrho$, тобто того мінімального
показника $\varrho$ з теореми Ia.
\end{zusatz}

Е.\ Бертіні\textsuperscript{10)} дав загальну формулу для $\varrho$.
Що ця формула, однак, не завжди дає мінімальне значення $\varrho$,
показано в цитованій праці на прикладі
$\varphi = x^{3} + y^{4}$, $\psi = x^{2} + y^{3}$. Нарешті зазначимо,
що й саму формулу Бертіні можна по-новому обгрунтувати нашою
головною теоремою.

\vspace{0.5em}
\noindent\rule{\textwidth}{0.4pt}
\vspace{0.3em}

{\footnotesize
\textsuperscript{7)} Якщо для $A(x,y)$ допустити також дробово-раціональні
функції $C(x,y):D(x,y)$, то $(A) = (C) - (D)$; отже в останньому
випадку $(A)$ може набувати і від'ємних значень.

\textsuperscript{8)} Пор. також доведення в зауваженні I ``Додатка''.

\textsuperscript{9)} Тут, отже, дістаються точні показники, на відміну від
оцінок панни Грете Герман [``Die Frage der endlich vielen Schritte
in der Theorie der Polynomideale'' (з використанням посмертних
теорем К. Гентцельта), \foreign{Math.\ Annalen} 95 (1926)], де, за
значно загальнішої постановки задачі, подаються лише верхні межі.

\textsuperscript{10)} E.\ Bertini, \foreign{Math.\ Annalen} 34 (1889).
}

\newpage
\section*{Теорема III і додаток}

Цікава спеціалізація головної теореми випливає з обмежувального
припущення, що нуль відповідного
$\mathfrak{q} = (\varphi, \psi, \mathfrak{p}^{e})$ є простою
точкою принаймні на одній із двох кривих $\varphi = 0$, $\psi = 0$.
Якщо, скажімо, він є простою точкою на $\psi$, то можна показати,
що число $t$ головної теореми точно дорівнює $\mu$, і що $t$
нерівностей головної теореми можна замінити однією-єдиною умовою
кратності. Це приводить до такої загальної теореми:

\begin{satz}[Теорема III]\label{satz:kap-III}
Якщо всі точки перетину $\varphi = 0$, $\psi = 0$ є простими
точками на $\psi$, то, за умови що $K$ взаємно простий з $\psi$,
для існування конгруенції $K \equiv 0(\varphi, \psi)$ необхідно й
достатньо, щоб $K$ зникав у всіх цих точках, причому так, щоб
кратність перетину в парі кривих $K = 0$, $\psi = 0$ у кожній з них
була принаймні такою великою, як у парі кривих
$\varphi = 0$, $\psi = 0$.
\end{satz}

Ця теорема III примітна ще з особливої причини: вимагана нею умова
кратності є такою, що її виконання або невиконання можна встановити
скінченним числом диференціювань, незалежно від формули
\eqref{eq:kap-2}, яка спирається на \eqref{eq:kap-1}. Це відразу
стає зрозуміло, якщо процитувати таку ``теорему'', яку я сформулював
1923 року\textsuperscript{11)}:

\smallskip
\textit{Якщо $A$ і $B$ --- будь-які два взаємно прості поліноми від
$x$, $y$, а $P$ --- їхня спільна точка, яка є простою точкою для
$B$, то кратність цієї точки як точки перетину тоді й лише тоді
дорівнює числу $\mu$, коли кратність тієї самої точки в парі кривих
$A' = 0$, $B = 0$ точно дорівнює $\mu - 1$; при цьому}
\[
A' = \frac{\partial A}{\partial x} \cdot \frac{\partial B}{\partial y} -
\frac{\partial A}{\partial y} \cdot \frac{\partial B}{\partial x}.
\]

Нарешті зазначимо ще, що головну теорему
(теорему~\ref{satz:kap-II}) з невеликою модифікацією можна перенести
також на однорідні поліноми від 3 змінних $x$, $y$, $z$. Це
перенесення в наведеному вище спеціальному випадку вдається особливо
елегантно. А саме можна показати:

\begin{zusatz}[до теореми III]
Теорема~\ref{satz:kap-III} залишається цілком незмінною, навіть за
формулюванням, якщо під $\varphi$, $\psi$, $K$ розуміти однорідні
поліноми від трьох змінних\textsuperscript{12)}.
\end{zusatz}

\vspace{0.5em}
\noindent\rule{\textwidth}{0.4pt}
\vspace{0.3em}

{\footnotesize
\textsuperscript{11)} H.\ Kapferer, ``Über die Multiplizität der Schnittpunkte von zwei algebraischen
Kurven''. \foreign{Jahresber.\ der D.\ M.-V.} 32 (1923).

\textsuperscript{12)} Питання про умови існування однорідної конгруенції від
трьох змінних $K(x,y,z) \equiv 0(\varphi(x,y,z), \psi(x,y,z))$,
хоч і за сильно спеціалізованих припущень, є також ядром дослідження
A.\ Ostrowski, ``Über die Maxwellsche Erzeugung der Kugelfunktionen'',
\foreign{Jahresb.\ der D.M.-V.} 33 (1925). У цій нотатці припускається:
$\varphi$ розкладається лише на лінійні множники; $\varphi$ і $K$
мають той самий порядок; $\psi = x^{2} + y^{2} + z^{2}$. Оскільки
тут $\psi$ цілком вільний від сингулярностей, головну теорему нотатки
Островського можна розглядати як спеціальний випадок нашої теореми III.
}

\newpage
\section*{Додаток, спільно з Е. Нетер\textsuperscript{13)}}

Система рівнянь \eqref{eq:kap-1} водночас дає раціональний метод
визначення повної системи лишків за $\mathfrak{q}$, під чим треба
розуміти повну систему представників скінченно багатьох, лінійно
незалежних відносно області коефіцієнтів $P$, класів лишків modulo
$\mathfrak{q}$. Завдяки цьому окремі кроки в доведенні головної
теореми стають дуже прозорими; водночас як побічний результат
дістаємо факт, що результантна кратність $\mu$ збігається з
довжиною ідеалу $\mathfrak{q}$;

\vspace{0.5em}
\noindent\rule{\textwidth}{0.4pt}
\vspace{0.3em}

{\footnotesize
\textsuperscript{13)} Ідеально-теоретичне тлумачення належить Е.\ Нетер, а
покладене в основу доведення (5) і (6) --- Г.\ Капфереру.

\textsuperscript{14)} Поле $P$ без обмеження загальності вважається
алгебраїчно замкненим; у цьому обсязі справді чинні всі попередні
теореми, за винятком умови диференціювання. Довжина $\mathfrak{q}$
визначається також як довжина композиційного ряду з ідеалів, що йде
від $\mathfrak{p}$ до $\mathfrak{q}$. Пор. E.\ Noether,
\foreign{Jahresber.\ d.\ d.\ Math.-Ver.} 34 (1925), S.~101
(коса пагінація); докладніше у H.\ Grell, \foreign{Math.\ Annalen} 97
(1927), S.~529.

Наявні в літературі доведення збігу обох чисел кратності (для $n$
змінних) дуже складні; найпростішим є ще не опубліковане посмертне
доведення Гентцельта.
}

\vspace{1em}
тобто з рангом кільця класів лишків або з числом лінійно незалежних
відносно $P$ класів лишків\textsuperscript{14)}.

Повна система лишків з позначеннями \eqref{eq:kap-1} і допоміжної
теореми II характеризується такою

\begin{satz}[Теорема IV]\label{satz:kap-IV}
Повна система лишків за $\mathfrak{q}$ задається поліномами
\[
h_{t}, h_{t}y, \ldots, h_{t}y^{(g_{t})-1}; \quad \ldots; \quad
h_{i}, h_{i}y, \ldots, h_{i}y^{(g_{i})-1}; \quad \ldots; \quad
h_{1}, h_{1}y, \ldots, h_{1}y^{(g_{1})-1}.
\]
При цьому кожного разу
$h_{i}, h_{i}y, \ldots, h_{i}y^{(g_{i})-1}$ дають повну систему
лишків від $\mathfrak{q}^{(i+1)}$ до $\mathfrak{q}^{(i)}$.
\end{satz}

З означення безпосередньо випливає:
$\mathfrak{q}^{(i+1)} = (\mathfrak{q}^{(i)}, h_{i})$, отже
$\mathfrak{q}^{(i)} \equiv 0(\mathfrak{q}^{(i+1)})$.
Теорему~\ref{satz:kap-IV} буде доведено, якщо показати:
\begin{equation}\label{eq:kap-4}
h_{i}x \equiv 0(\mathfrak{q}^{(i)}); \qquad
h_{i}y^{(g_{i})} \equiv 0(\mathfrak{q}^{(i)})
\end{equation}
і
\begin{equation}\label{eq:kap-5}
\text{з } h_{i}F(y) \equiv 0(\mathfrak{q}^{(i)})
\quad\text{випливає}\quad F(y) \equiv 0(y^{(g_{i})}).
\end{equation}

Співвідношення \eqref{eq:kap-4} саме й кажуть, що лінійні комбінації
$h_{i}, h_{i}y, \ldots, h_{i}y^{(g_{i})-1}$ вичерпують класи лишків
від $\mathfrak{q}^{(i+1)}$ до $\mathfrak{q}^{(i)}$, тоді як
\eqref{eq:kap-5} виражає лінійну незалежність. Оскільки
$\mathfrak{q}^{(t+1)}$ стає одиничним ідеалом, що тепер випливає зі
скінченної довжини $\mathfrak{q}$, звідси одержується доведення
теореми~\ref{satz:kap-IV} і водночас факт, що довжина
$\mathfrak{q}$ дорівнює $(g_{1}) + \cdots + (g_{t})$, тобто за
\eqref{eq:kap-2} збігається з результантною кратністю.

Перше зі співвідношень \eqref{eq:kap-4} читається з
\eqref{eq:kap-1}; друге дістається так:

Маємо $h_{i}g_{i} \equiv 0(\mathfrak{q}^{(i)})$; отже, через перше
співвідношення \eqref{eq:kap-4}, також
$h_{i}g_{i}(0,y) \equiv 0(\mathfrak{q}^{(i)})$, або
$h_{i}y^{(g_{i})}(g_{i}(0,y):y^{(g_{i})})
\equiv 0(\mathfrak{q}^{(i)})$, звідки
$h_{i}y^{(g_{i})} \equiv 0(\mathfrak{q}^{(i)})$ через
$g_{i}(0,y):y^{(g_{i})} \not\equiv 0(y)$.

Перед доведенням \eqref{eq:kap-5} зауважимо те, що безпосередньо
зчитується з \eqref{eq:kap-1}: $f_{i}$ і $g_{i}$ можуть мати
спільним дільником лише поліном $d_{i}(y)$ від $y$, не подільний на
$y$: $f_{i} = d_{i}(y)\bar{f}_{i}$,
$g_{i} = d_{i}(y)\bar{g}_{i}$, причому $\bar{f}_{i}$ і
$\bar{g}_{i}$ взаємно прості. Оскільки
$d_{i}(y) \not\equiv 0(\mathfrak{p})$, маємо також
$\mathfrak{q}^{(i)} = (\bar{f}_{i}, \bar{g}_{i}, \mathfrak{p}^{e})$;
при цьому це дорівнює
$(\bar{f}_{i}, \bar{g}_{i}, \mathfrak{p}^{\sigma})$ для кожного
$\sigma \geq \varrho$, як дільник $\mathfrak{q}$ для кожного такого
$\sigma$; отже $\mathfrak{q}^{(i)}$ є первинною компонентою
$(\bar{f}_{i}, \bar{g}_{i})$, що відповідає нулю.

Нехай $\gamma_{1}, \delta_{1}; \ldots; \gamma_{r}, \delta_{r}$ ---
нулі $(\bar{f}_{i}, \bar{g}_{i})$, відмінні від $x = 0$, $y = 0$, і
нехай $\tau_{1}, \ldots, \tau_{r}$ --- показники відповідних
первинних компонент $(\bar{f}_{i}, \bar{g}_{i})$; далі в добутку
\[
P = ((x-\gamma_{1}) + u(y-\delta_{1}))^{\tau_{1}} \cdots
((x-\gamma_{r}) + u(y-\delta_{r}))^{\tau_{r}}
\]
невизначену $u$ виберімо як елемент з $P$ так, щоб збереглося
$P \not\equiv 0(\mathfrak{p})$. Якщо покласти $G(x,y)$ рівним
добутку спеціалізованого $P$ з $d_{i}(y)$, то
$G(x,y) \not\equiv 0(\mathfrak{p})$ і
$\mathfrak{q}^{(i)} \cdot G(x,y) \equiv 0(f_{i}, g_{i})$.

З припущення $h_{i}F(y) \equiv 0(\mathfrak{q}^{(i)})$ випливає
$h_{i}F(y)G(x,y) \equiv 0(f_{i}, g_{i})$, або, разом з
\eqref{eq:kap-1}, у формі тотожностей:
\[
h_{i}F(y)G = A f_{i} + B g_{i}; \qquad
h_{i}x = (g_{i}(0,y):y^{(g_{i})})f_{i}
       - (f_{i}(0,y):y^{(g_{i})})g_{i}.
\]
А це дає
\[
\bar{f}_{i}(xA - F(y)H) \equiv 0(\bar{g}_{i}); \qquad
H = G(x,y) \cdot g_{i}(0,y):y^{(g_{i})}
\not\equiv 0(\mathfrak{p}).
\]
З взаємної простоти $\bar{f}_{i}$ і $\bar{g}_{i}$ далі випливає
\[
xA - F(y)H \equiv 0(\bar{g}_{i});
\]
отже $F(y)H \equiv 0(\bar{g}_{i}, x)
\equiv 0(\bar{g}_{i}, x, \mathfrak{p}^{\sigma})
= (y^{(g_{i})}, x)$, а звідси $F(y) \equiv 0(y^{(g_{i})})$.
Тим самим доведено \eqref{eq:kap-5}, а отже й
теорему~\ref{satz:kap-IV}.

\smallskip
\textbf{Зауваження I.} Співвідношення \eqref{eq:kap-4} і
\eqref{eq:kap-5} разом із допоміжною теоремою II можна записати як
взаємно обернені співвідношення між частками ідеалів\textsuperscript{15)};
при цьому доведення другого співвідношення водночас дає доведення
допоміжної теореми II:
\begin{equation}\label{eq:kap-6}
\mathfrak{q}^{(i)} : \mathfrak{q}^{(i+1)}
= (\mathfrak{q}^{(i)}, x) = (y^{(g_{i})}, x); \qquad
\mathfrak{q}^{(i)} : (\mathfrak{q}^{(i)}, x)
= \mathfrak{q}^{(i+1)}.\text{\textsuperscript{16)}}
\end{equation}
Перше співвідношення є підсумуванням \eqref{eq:kap-4} і
\eqref{eq:kap-5}; щодо другого треба зауважити: у
\eqref{eq:kap-4} було доведено
$x \cdot \mathfrak{q}^{(i+1)} \equiv 0(\mathfrak{q}^{(i)})$;
обернене дістається аналогічно до \eqref{eq:kap-5}. З
$xM \equiv 0(\mathfrak{q}^{(i)})$ як вище випливає:
\[
xMG \equiv 0(f_{i}, g_{i}); \quad\text{отже}\quad
xMG \cdot g_{i}(0,y):y^{(g_{i})} \equiv 0(xh_{i}, g_{i}).
\]
Оскільки $g_{i} \not\equiv 0(x)$ за означенням, а саме
$(f_{i}) \geq (g_{i})$, звідси виходить
\[
M \cdot (G \cdot g_{i}(0,y):y^{(g_{i})})
\equiv 0(h_{i}, g_{i}) \equiv 0(\mathfrak{q}^{(i+1)}),
\quad\text{тобто}\quad M \equiv 0(\mathfrak{q}^{(i+1)}).
\]

Необхідність умови допоміжної теореми II тепер випливає так:
$K_{i} \equiv 0(\mathfrak{q}^{(i)})$ дає
$K_{i}(0,y) \equiv 0(y^{(g_{i})}, x)$; отже
$(K_{i}) \geq (g_{i})$, а тим самим, за \eqref{eq:kap-3},
\[
x \cdot K_{i+1} \equiv 0(\mathfrak{q}^{(i)});
\quad\text{отже, за другим співвідношенням \eqref{eq:kap-6},}\quad
K_{i+1} \equiv 0(\mathfrak{q}^{(i+1)}).
\]

Достатність умови зчитується безпосередньо; скінченне повторне
застосування допоміжної теореми II дає головну теорему. Так само
ітерація другого співвідношення \eqref{eq:kap-6} дає додаток I, а
саме $\mathfrak{q} : x^{t} = \mathfrak{q}^{(t+1)}$, тобто одиничний
ідеал; причому $t$ є найнижчим таким степенем, оскільки
$\mathfrak{q} : x^{t-1} = \mathfrak{q}^{(t)}$.

\textbf{Зауваження II.} Для полінома $K$, не подільного на
$\mathfrak{q}$, можна ще сказати: якщо
$K \equiv 0(\mathfrak{q}^{(i+1)})$,
$\not\equiv 0(\mathfrak{q}^{(i)})$, то
$(K, \mathfrak{q}^{(i)}) = (h_{i}y^{(K)}, \mathfrak{q}^{(i)})$.
Бо за теоремою~\ref{satz:kap-IV}
$K \equiv h_{i}y^{\lambda}c(y)$ modulo
$(\mathfrak{q}^{(i)})$ з $c(y) \not\equiv 0(y)$; отже
\[
(K, \mathfrak{q}^{(i)}) = (h_{i}y^{\lambda}, \mathfrak{q}^{(i)}).
\]
При цьому $y^{(g_{i})-\lambda}K \equiv 0(\mathfrak{q}^{(i)})$, і саме це
за \eqref{eq:kap-5} є найнижчим таким степенем. За допоміжною
теоремою II цей найнижчий степінь дорівнює $(g_{i}) - (K)$; отже
$\lambda$ дорівнює $(K)$.

\vfill
\begin{flushright}
(Надійшло 24.\ 10.\ 1926.)
\end{flushright}

\vspace{0.5em}
\noindent\rule{\textwidth}{0.4pt}
\vspace{0.3em}

{\footnotesize
\textsuperscript{15)} Під часткою $\mathfrak{a}:\mathfrak{b}$, як відомо,
розуміють сукупність поліномів $c(x,y)$ таких, що
$\mathfrak{b}c \equiv 0(\mathfrak{a})$.

\textsuperscript{16)} Із цих двох взаємно обернених співвідношень одне
зумовлює інше за загальною теорією незвідних ідеалів; пор.
Macaulay, \textit{Modular Systems}, Nr.\,72 і 73.
\foreign{Cambridge Tracts} 19 (1916).
}

\end{document}
